\typesetchapter{Chapter 3}{Mountain}

\noindent Sa-Kailan did not sleep well after the first lesson.\par

His body wanted sleep. His mind would not close. It kept returning to the basin, to the token at the bottom, to the sign of Light shifting under water. The thing was small enough to be embarrassing. A child could drop a spoon into a cup and see it bend. A worker could rinse his hands in a washing pool and watch his fingers become pale and broken beneath the surface. Every person in the Mountain had seen water make the familiar untrustworthy.\par

But Sa-Kailan had not merely seen it. He had noticed it.\par

That was different. Not understanding. Not yet. Understanding belonged to men like Va-Raedin, who could ruin a boy’s pride with one question and call it instruction. But the sight had lodged inside him. Light entered water and came out changed. It crossed a threshold and lost its obedience. It could be delayed, diverted, thickened, misplaced. The basin by his mat was no longer only a basin. The lamp above it was no longer only a lamp. The faint wavering on the wall where a servant had spilled water was no longer accident. Everything seemed to be doing something behind the name by which he had known it.\par

By morning he felt both exhausted and sharpened. He dressed before the household fully woke, tied his robe with more care than usual, and checked the bandage around his thumb. The cut had closed badly during sleep. His mother had left dried fruit and a small clay cup of goat’s milk in the wall recess. The milk had formed a skin. He moved the cup under the lamp and watched the yellow surface tremble as his fingers touched the rim.\par

Then hunger remembered itself, and he drank.\par

The day after an Ascension always moved unevenly. The Mountain had spent order on ceremony and now had to earn it back through work. Those who had stood too long in the Hall still had to carry, cook, count, repair, wash, teach, bless, guard, and record. The old had been kept warm through the rite and now required broth and washing. Children who had slept through the closing prayer had woken late and irritable. Faith had new instructions to distribute under Ka-Elun, Keeper of Faith. Justice had the transfer to reconcile with old tablets. Light had a new Candidate and expected him before ordinary duties. That meant someone else had been moved from some other task. That meant Balance would have to adjust a ledger.\par

Balance always noticed.\par

He did not get far. At the first turn, a Shield warden raised one hand and stopped three workers, two children, and Sa-Kailan together.\par

“South service way is closed,” the warden said.\par

One of the workers, a woman carrying folded cloth bundles, frowned.\par

“Since when?”\par

“Since it was closed.”\par

“I have cloth due in the herb rooms.”\par

“Then it will arrive by another passage.”\par

“That adds half a descent.”\par

The warden looked at the bundles, not the woman.\par

“Then the cloth will be tired.”\par

The children laughed before their father quieted them. The worker did not. She shifted the bundles against her hip, irritated but not surprised. No one asked why the passage had closed. In the Mountain, closure did not require explanation from Shield. It required adjustment from everyone else.\par

The warden looked at Sa-Kailan’s robe and bandaged thumb.\par

“Light?”\par

“Candidate,” Sa-Kailan said.\par

The warden’s gaze sharpened by a fraction.\par

“Sa-Kailan?”\par

The new prefix still entered him strangely.\par

“Yes.”\par

“You are expected east. Use the lower wash route. Do not cross the copper gate. It is being witnessed.”\par

Being witnessed. Justice, then.\par

Sa-Kailan nodded and turned with the others. Shield did not explain closures. Justice did not explain witness unless sequence demanded it. Light might have found a way to see without asking. He glanced back once as the warden moved aside for a woman in Balance green carrying a sealed slate. She passed through the closed way without question.\par

Permission, Sa-Kailan thought, was not a wall. It was a sorting mechanism.\par

The lower wash route took him through the waking parts of the Mountain. Here the stone was less polished than in Light. The passages were narrower, warmer, and more honest in their labour. A boy dragged a basket of cracked jars toward the repair rooms. Two old women argued softly over a cloth filter. A man with flour on his wrists slept upright against a wall while waiting for a gate to release. Above them, water knocked faintly inside hidden pipes, not loud enough to disturb, not quiet enough to forget.\par

The first water queue had already formed along a shallow channel under the supervision of two Balance clerks and one Faith attendant. Each family had a marked vessel. Some were clay, some metal, a few old ones carved from dark stone and repaired so often that the repairs had become part of their identity. A copper rod lay across the channel. The clerks moved it from notch to notch, reading allocation by household mark, work assignment, infirmary allowance, pregnancy, heat exposure, and yesterday’s draw.\par

The Faith attendant blessed each jar before the measure was released. Her prayer did not slow the work. It timed it.\par

\begin{scripture}
“God who hid the waters beneath the Mountain, teach us to count without cruelty,” she said over one jar. The Balance clerk did not look up.
\end{scripture}

“Half measure.”\par

The woman holding the jar stiffened.\par

“Yesterday was half measure.”\par

“Yesterday was second descent. Today is third. Your household drew early after the Ascension heat.”\par

“My father is ill.”\par

“Your father was ill yesterday. Illness is recorded. Extra warmth was granted. Water was advanced.”\par

The Faith attendant touched the woman’s wrist before anger could become public.\par

“The hidden waters remember the unborn also.”\par

The woman’s mouth tightened. Briefly she seemed about to answer. Then she lowered her eyes and took the half-filled jar.\par

Sa-Kailan watched the water settle inside it. The surface caught a high lamp and bent the light along the curved clay. The jar looked fuller from one angle than another. He knew it was half measure because Balance had named it half measure. Without the mark, the eye might have argued.\par

“Candidate of Light.”\par

It was not a greeting.\par

“Yes?”\par

“You are blocking the return line.”\par

Sa-Kailan moved at once, embarrassed. The clerk’s face remained neutral, which made the correction more effective. Balance had no need to scold where inconvenience could be named precisely.\par

Beyond the water queue, the kitchens opened in a series of low heated chambers. The air changed again. It became thick with steam, yeast, boiled grain, goat milk, dried herbs, smoke from carefully managed dung fuel, and the sour-sweet smell of fermentation. The kitchen workers were mostly Balance, though Faith was everywhere in their timing: short prayers for sealing jars, longer ones for opening stored grain, work chants that kept hands moving together when the heat made speech unpleasant.\par

“Let every fruit be gathered. Let every vessel be sealed. Let no waste enter the Oasis.”\par

The workers did not sing beautifully. Some did not sing at all. But the words entered the labour without needing to be believed in the same way by every mouth. One woman sang with devotion. Another used the lines to count kneading turns. A boy mouthed them only when an elder looked his way. Faith, Sa-Kailan realized, did not require equal sincerity from every voice. It required the rhythm to hold.\par

Two Balance girls dusted trays with powdered shell. Calcium, he remembered from childhood lessons. Shell belonged to no romance, no prayer, no myth. It belonged to bones, teeth, poultice, mortar, and the quiet insistence of bodies.\par

At the far end of the kitchens, three goats were painted on a stone panel above a locked hatch. The animals looked almost cheerful, with square beards and solemn eyes. Beneath them, a line had been carved in Balance script and copied below in common hand:\par

THE OUTSIDE FEEDS ONLY WHAT THE MOUNTAIN CAN COUNT.\par

Sa-Kailan slowed.\par

Outside. The word struck differently now. The Oasis. The shadow of forest in the old teachings. The permitted prayer periods. Balance workers who tended animals and plants beyond stone. Shield wardens who watched paths no one else was supposed to understand. Children were taken to pray before the Oasis at appointed ages, but always in groups, always under instruction, always with the eyes of Shield behind them and the voice of Faith before them. The Oasis was water, food, memory, boundary, temptation. It was the world God had left them, or the fragment of world the Mountain had not yet consumed.\par

A kitchen supervisor saw him staring.\par

“You want goat duty, Light?”\par

A few workers laughed.\par

“No,” Sa-Kailan said.\par

“Then keep moving. Even candidates cast shadows.”\par

He moved.\par

The washing pools lay below the kitchens, not for dignity but for sequence. Water that had rinsed vegetables and cooled sealed jars was filtered, blessed, measured again, and used where it could be used without insulting health. Balance called it sequence. Faith called it gratitude. Children called it old water and were corrected until they stopped. Nothing in the Mountain ended after one use if it could be made to serve again.\par

Near the tool wash, Sa-Kailan stopped. A woman dipped a curved metal blade into a shallow pool. The blade bent under the surface. Not truly. He knew that. But the image turned at the threshold between air and water. The woman scrubbed it with sand, lifted it out, and the blade returned to itself.\par

Sa-Kailan felt again the sudden brightness from the night before.\par

Water made the familiar untrustworthy. Not by hiding it, but by changing the path by which sight reached it. And everyone walked past this all day. Workers, clerks, singers, wardens. Even Light attendants, perhaps. They named it ordinary and so it remained hidden.\par

“Sa-Kailan?”\par

He turned too quickly.\par

Sa-Tavan stood at the edge of the wash passage with a slate under one arm and the new black cord visible at his wrist. He looked as he had in the Hall: narrow, composed, and faintly disapproving, though perhaps that was only the face Justice gave its children before they earned a better one.\par

“You are far from the east passage,” Sa-Tavan said.\par

“So are you.”\par

“I have a permission mark.”\par

He lifted the slate. A small piece of pale wax had been pressed onto its upper corner and marked with the sign of Justice.\par

“Congratulations.”\par

“It is not an honour. It is a route allowance.”\par

“That was not what I meant.”\par

“It was what you said.”\par

Sa-Kailan almost smiled despite himself.\par

“You were less exact yesterday.”\par

“I was being named yesterday.”\par

“Does Justice permit inexactness during ceremony?”\par

Sa-Tavan considered this with enough seriousness that Sa-Kailan regretted the joke.\par

“Ceremony permits formal compression. It does not permit falsehood.”\par

“Useful distinction.”\par

“Yes.”\par

The answer was so earnest that it became impossible to mock. Sa-Kailan looked at the slate.\par

“What are you recording?”\par

“Disputed sequence in the second water queue.”\par

“What dispute?”\par

“A household claims the mark was advanced after the jar was filled. Balance claims the mark was advanced before. Faith witnessed the blessing but not the measure. So Justice records the order while memory is still warm.”\par

Sa-Kailan looked back toward the water lines. The woman with the half measure was gone. Others had taken her place.\par

“And if you record it correctly?”\par

“Then the next measure can be corrected if correction is required.”\par

“And if you record it incorrectly?”\par

Sa-Tavan’s face changed slightly. Not fear. Respect.\par

“Then a correction becomes a wound.”\par

It was an answer Sa-Kailan would not have expected from him. He had thought Justice dry, concerned with tablets and cords and the satisfaction of being right after everyone else had grown tired. But the boy’s seriousness had weight beneath it. A measure misrecorded became hunger. A witness misplaced became punishment. A sequence altered became guilt where none belonged.\par

Sa-Kailan nodded.\par

“Then I should not delay you.”\par

“No.”\par

Sa-Tavan turned to go, then paused.\par

“Your thumb is bleeding through.”\par

Sa-Kailan looked down. A small red mark had appeared on the cloth.\par

“It opened during training.”\par

“You should have it witnessed in the infirmary if Light cut it ceremonially and training reopened it.”\par

“Why?”\par

“So responsibility is not confused.”\par

“Between my thumb and Light?”\par

“Between ritual wound and instruction wound.”\par

Sa-Kailan stared at him, then laughed once. Sa-Tavan did not laugh.\par

“If it scars badly, the record should know why.”\par

He left before Sa-Kailan could answer.\par

For a few moments, Sa-Kailan stood among the steam and washing noise, irritated and impressed in equal measure. He had been given a title and thought it made him more visible. Sa-Tavan had been given a cord and already saw the world as a chain of possible misplacements. It was absurd. It was also useful. Justice waited not because it loved delay, but because action without sequence could change the meaning of everything that followed.\par

He decided not to ask anyone to witness his thumb.\par

The child halls were higher, carved into warmer stone near the family districts and the teaching chambers of Faith. Sa-Kailan passed them on his way toward the eastern lamp corridors, though again he took the long route. He told himself he was observing. That made the delay sound like discipline rather than avoidance.\par

Children sat in three groups on woven mats while a Faith teacher moved between them with a cord of five knots. The youngest repeated after her in singsong voices, stumbling over the old forms.\par

“The Mountain is refuge.”\par

“The Mountain is refuge.”\par

“The Oasis is water.”\par

“The Oasis is water.”\par

“The forest was plenty.”\par

“The forest was plenty.”\par

“Beyond the forest thou shalt not wander.”\par

“Beyond the forest thou shalt not wander.”\par

The teacher paused and touched the fifth knot.\par

“And if the forest is gone?”\par

The youngest children faltered. The older ones answered for them.\par

“The Covenant remains.”\par

Sa-Kailan stopped near the doorway out of habit rather than intent. Yesterday the words would have sounded complete. Now they sounded like the surface of something harder. The forest was gone. The Covenant remained. If a child did not know what had been lost, the line comforted. If an adult knew, it instructed grief to behave.\par

On the far wall of the teaching chamber hung five small panels for the Kahir. Balance was represented by a cup and sealed jar. Faith by an open text. Shield by a gate. Light by an aperture. Justice by a black cord crossing a pale stone. The children were too young to be afraid of them. They pointed and whispered and tried to guess which office would one day claim them, as if choosing sweets from a festival tray.\par

One girl pointed to Light and said, “They know all the secret doors.”\par

The teacher turned.\par

“No.”\par

The girl shrank.\par

“Light knows that doors are not the first secret,” the teacher said.\par

Sa-Kailan almost stayed to hear the rest.\par

He moved on.\par

The lamp corridors lay along the eastern inner spine. Here Light became less ceremonial and more necessary. The Mountain could not live in darkness, but it could not flood itself with flame either. Oil had to be rationed. Smoke had to be drawn. Mirrors had to be cleaned without scratching. Lamps had to serve three purposes at once: guide workers, reveal danger, and conceal routes not meant for every eye. Narrow inspection slots allowed Light attendants to see whether distant lamps burned correctly without entering passages where they had no right to stand.\par

A senior candidate of Light stood on a ladder adjusting a small mirror while two apprentices held the base steady. The mirror was no larger than a plate, but it directed light fifty paces down a side corridor. The beam struck a patch of pale stone near a stair, then broke into a softer glow that made the steps visible without lighting the sealed door beyond them.\par

“Too high,” the candidate said.\par

One apprentice touched the bracket.\par

“Not with your hand. With the screw.”\par

“The screw is stuck.”\par

“The screw is not stuck. Your patience is.”\par

Sa-Kailan smiled despite himself.\par

The candidate saw him.\par

“Fellow Candidate. Since you are smiling, tell me what is wrong with this mirror.”\par

The apprentices turned with obvious relief at someone else being chosen. Sa-Kailan studied the mirror, the lamp, the side corridor, the stair. The obvious answer was that the beam sat too high. She had already said so, which meant that answer would embarrass him.\par

He walked to the side, then crouched. From below, he saw the faint duplicate of the beam along the lower wall. A small film of dust had gathered at the mirror’s bottom edge, scattering enough light to touch the outline of the sealed door.\par

“It is not only too high,” he said. “It is dirty.”\par

The candidate looked at him.\par

“Everything is dirty.”\par

“The dirt is revealing the door.”\par

“Better.”\par

Not Va-Raedin’s voice, not Va-Raedin’s approval, but the word still struck. She handed one apprentice a cloth.\par

“Clean the lower edge. Then lower by two turns, not three. Three will light the door frame. Two will light the stair.”\par

Sa-Kailan moved on before she found another question.\par

He liked the lamp corridors more than he wanted to. They showed Light in its least pious form. No chanting. No gravures. No blue-rooted ceremonial flames bending toward holy faces. Here Light was labour: oil, soot, wick, screw, angle, dust, air. It was infrastructure disguised as grace, or grace made possible by infrastructure. He could not decide which thought made him more uncomfortable.\par

At the repair spaces near the lower spine, men and women worked with cracked mirror plates, bent brackets, broken hinges, jar seals, cart wheels, pump handles, and pipe sections. The repair halls belonged to no Kahir entirely. Balance inspected materials. Light inspected reflective surfaces. Shield inspected locks and hinges. Justice witnessed anything whose failure had caused injury. Faith blessed new work when the repair touched a sacred route. The result was a constant low-level argument restrained by procedure.\par

“It will hold,” a Balance man said, turning a pipe section in both hands.\par

A Light woman shook her head.\par

“It reflects into the side slit.”\par

“It is pipe. It is not meant to reflect.”\par

“It reflects whether you intend it or not.”\par

“Then dull it.”\par

“If I dull it, it corrodes faster.”\par

Justice, in the person of an elderly recorder seated with a slate on her lap, said without looking up, “State the consequence of each choice.”\par

The Balance man sighed.\par

“If polished, it misdirects light.”\par

The Light woman said, “If dulled, it weakens sooner.”\par

“Time to failure?”\par

“Unknown.”\par

“Then it is not yet a choice,” the recorder said. “It is two fears.”\par

Sa-Kailan slowed. The recorder’s voice had no force in it, yet both workers stopped arguing. That was Justice too, then. Not only judgment after harm, but the discipline of forcing a disagreement to name itself before it became a quarrel.\par

The medical and herb rooms lay beyond the repair halls, under warmer vents where fungus trays and medicinal roots could be kept at their proper dampness. The air there was rich, green, and bitter. Bundles of herbs hung from cords. Clay trays held moss, pale fungi, and small red leaves grown under cloth-filtered lamps. Mortars thudded. A child with a burn on his forearm sat very still while a healer applied paste. A Faith attendant murmured comfort. A Balance clerk noted the amount of paste, oil, cloth, and future dressing the child would require.\par

“Will it scar?” the mother asked.\par

The healer did not answer at once. She wrapped the arm carefully, then glanced toward the recorder.\par

“It will mark,” she said.\par

The Faith attendant touched the child’s hair.\par

“A mark is not always a wound.”\par

The Balance clerk said, “He will need extra oil in the dressing for three days.”\par

The mother looked frightened by the word extra.\par

The clerk added, “Recorded under injury allowance.”\par

Relief moved through the woman so visibly that Sa-Kailan felt embarrassed for having witnessed it. Mercy, here, was not a feeling. It was a category. If the burn was recorded correctly, the child received oil. If not, pity had to fight the storehouse, and pity did not always win.\par

On the far side of the herb room, stacked wooden frames held fibre plants under careful light. Sa-Kailan recognized flax, nettle, and two other plants he knew by use but not name. Cloth, cord, bandage, wick. Everything connected. Goats outside gave milk, hide, sinew. Shell gave powder. Bees gave wax and honey. Herbs gave medicine. Fibre gave binding. Waste gave fuel, feed, compost, or warning. Balance did not simply count the Mountain’s life. It arranged it so that nothing ended at its first purpose.\par

By midday, Sa-Kailan had reached the storehouses. He had not meant to. Light candidates had no business lingering there without assignment, but the long route had become a private map of the Mountain’s hidden relations, and the storehouses were where relation became visible.\par

They lay behind three gates. The first stood open and busy. The second required a mark. The third was watched by Shield and witnessed by Justice. Sa-Kailan remained in the outer hall, where permitted goods were received, sorted, inspected, and redirected.\par

The storehouse air was dry and cooler than the kitchens. It smelled of grain, wax, old wood, salt, dried fruit, shell powder, leather, and the faint sharpness of preserved herbs. Rows of jars stood against the walls, each marked by date, household source, use category, and restriction level. Balance workers moved among them without hurry. Haste wasted attention. Attention preserved life.\par

At a central table, a woman in Balance green inspected a tray of dried figs. She cut one open, frowned, and set it aside. Another. Another. The worker who had brought the tray shifted nervously.\par

“Mould?” he asked.\par

“Beginning.”\par

“How much lost?”\par

“That depends on whether you argue.”\par

He closed his mouth.\par

She separated the figs into three piles: clean, questionable, spoiled. The spoiled went into a sealed waste pot. The questionable were marked for boiling. The clean were reweighed. A Faith attendant beside her spoke a line over the clean figs, then a different line over the waste pot. Sa-Kailan had never noticed that before. Even waste received language. Not the same blessing, but acknowledgment. Nothing left the order of the Mountain without being named.\par

“Sa-Kailan.”\par

This time the voice was older, female, and carried authority enough that he straightened before turning.\par

Ka-Leth, Keeper of Balance, stood at the entrance to the second gate. He had seen her in ceremony, broad and grave, pouring water with hands stained faintly green. Up close she seemed less ceremonial and more dangerous. She looked at his robe, his bandage, his face, and the distance between his feet and the restricted line painted across the floor.\par

“Are you assigned here?”\par

“No, Keeper.”\par

“Are you lost?”\par

“No, Keeper.”\par

“Then you are either idle or observing.”\par

“Observing.”\par

“Light’s favourite excuse.”\par

Sa-Kailan lowered his eyes.\par

“I can leave.”\par

“You can. But since you have observed, tell me what you saw.”\par

He looked at the figs, the jars, the marks, the waste pot, the gates, the clerk’s blade, the Faith attendant’s mouth, the line on the floor he had nearly crossed. He did not want to give a childish answer. He also did not want to stand silent before Ka-Leth until silence became confession.\par

“Nothing is only what it first appears to be,” he said.\par

Ka-Leth’s face did not change.\par

“That is a Light answer. Try again.”\par

He swallowed.\par

“The spoiled figs are not simply thrown away. They are sealed. The questionable ones are not saved as they are, but changed into something safer. The clean ones are weighed again because loss changes the measure. The prayer changes with use. Waste is still counted.”\par

“Better. Why?”\par

“Because waste that is not counted becomes invisible.”\par

“And invisible waste becomes?”\par

He thought of the water queue. The injured child. The pipe in the repair hall. The half-filled jar that looked fuller from the wrong angle.\par

“Theft from someone who does not know he has been robbed.”\par

Ka-Leth considered him.\par

“Va-Raedin is not wasting you entirely.”\par

The praise was so severe it almost passed as insult.\par

“Thank you, Keeper.”\par

“Do not thank me. You are too near the second gate. Step back.”\par

He stepped back.\par

Ka-Leth turned away, then paused.\par

“If Light wishes to understand the Mountain, it should remember that revelation consumes oil. Curiosity is not free.”\par

She went through the gate, and Shield closed it behind her.\par

Sa-Kailan stood there longer than was wise, the sentence settling. Curiosity is not free. From Va-Raedin, it would have meant pride costs humiliation. From Balance, it meant something colder. Time, oil, warmth, water, attention, risk. Every question required a lamp, a corridor, a witness, a pause in someone else’s work. Even discovery drew from the storehouse.\par

By the time he returned toward the east, the Mountain had entered its afternoon rhythm. Early hunger had passed into work. Children were in lessons or assigned tasks. The kitchens had shifted from preparation to cleaning. Water queues shortened and lengthened according to district. Shield changed watches at inner thresholds with no bell and little speech. Justice tablets lay open near three gates, weighted by smooth black stones while ink dried.\par

From a side passage he heard Shield voices, low and formal.\par

“Gate held.”\par

“Gate witnessed.”\par

“Gate held.”\par

“Gate remembered.”\par

Sa-Kailan slowed but did not turn toward them. Shield would appear in full soon enough. At an inner gate perhaps, or near one of the outside routes. Shield always watched the place where permission became exposure.\par

The lamp above the eastern return arch had been turned. Only a little. Enough that the polished strip along the wall sent light downward instead of across. The floor near the arch showed faint wetness where a servant had mopped and not fully dried the stone. The reflected lamp trembled in the water.\par

He crouched.\par

The servant saw him and stopped, cloth in hand.\par

“Candidate?”\par

He stood too quickly.\par

“Nothing.”\par

She looked at the floor, then at him.\par

“It was only wash water.”\par

Only wash water.\par

He almost laughed. Not at her. At the Mountain, perhaps. At himself. At the arrogance of thinking secrets announced themselves as secret.\par

“Yes,” he said. “I know.”\par

The servant waited, unsure whether Light expected something further. When he said nothing, she walked on.\par

Sa-Kailan continued east with the sense that the ordinary world had become densely written. Not in words. In relations. Water and light. Prayer and valve. Gate and witness. Food and permission. Injury and oil. Waste and theft. A child’s burn and a storehouse category. A closed passage and a woman in green carrying a slate through it. Every office claimed only one portion of the whole, but the whole existed through their collisions.\par

That realization should have comforted him. It was The Covenant made practical. Let no one hand hold the whole.\par

But the same realization tightened around him. If no one hand held the whole, then every person lived inside arrangements he could not see. Every blessing might conceal a valve. Every closed gate might conceal danger or merely fear. Every measure might be mercy or theft depending on the record behind it. Every silence might protect him, or keep him obedient to someone else’s protection.\par

Near the eastern stair, he passed the entrance to a narrow prayer alcove. Two old men sat inside, heads bowed before a small bowl of water and a lamp no larger than a fist. Between them, a child recited from memory while a Faith attendant corrected his rhythm.\par

\begin{scripture}
“And God said, Keep thou the Mountain, and keep thou the Oasis, and keep thou the Temple.” The child stumbled. The attendant waited.
\end{scripture}

“Keep the secrets apart,” the child continued, “and let no one hand hold the whole.”\par

Sa-Kailan stopped just beyond the alcove, where they could not see him without turning. He had said those words in the Hall. Now they followed him differently. Keep the secrets apart. Not as poetry. As design. As law. As threat.\par

He returned at last to the candidate quarters and found them half empty. Some had already gone to work assignments. Others slept in the careless positions of youth that had not been trained out of them yet. Sa-Kailan sat on the edge of his sleeping mat and removed the Light token from the pouch at his belt.\par

He held it up toward the lamp.\par

The token did not show him more. It showed him less with greater force. Light did not always reveal by widening sight. Sometimes it revealed by making the eye admit it could look only through one opening at once.\par

Outside the chamber, the Mountain continued without him. Water moved. Bread cooled. Gates were held. Records dried. Children recited. Lamps burned by measure, not by desire. Somewhere, Balance counted what Faith blessed. Somewhere, Justice wrote what Shield would not explain. Somewhere, Va-Raedin knew which mirror he had failed to notice. Somewhere, Va-Sheva stood at a threshold he was not permitted to approach.\par

Sa-Kailan closed his fingers around the token. He had been proud of becoming Candidate of Light. He still was. The pride had not diminished; if anything, the day had fed it. But pride alone was no longer clean enough. He had wanted Light to lift him above the ordinary machinery of the Mountain. Instead Light had made the machinery visible. Not all of it. Just enough.\par

The Mountain was refuge. The Mountain was order. The Mountain was the reason the hidden people had not vanished into the open world and its cruelties. He knew this. He believed it, or believed that he believed it. The Mountain was also a prison made of permissions. Water, bread, light, silence, warmth, prayer, passage, record, touch, even the right to stand before the Oasis and look upon the world God had taken from them and promised one day to return — all of it moved through hidden structures. Nothing was only itself. Nothing arrived unshaped.\par

Sa-Kailan opened his hand. The token lay there, small and heavy, marked with the sign of the Kahir that had claimed him. He looked through its aperture one last time at the basin beside his mat. The water inside was still, but he no longer trusted stillness. He had begun to see the Mountain. And the Mountain, he understood with a discomfort close to awe, had been seeing him all along.\par

\ornament
